\documentclass[11pt, a4paper]{article}

% ── Packages ──────────────────────────────────────────────────────────
\usepackage[utf8]{inputenc}
\usepackage[T1]{fontenc}
\usepackage{lmodern}
\usepackage[margin=1in]{geometry}
\usepackage{amsmath, amssymb, amsthm}
\usepackage{booktabs}
\usepackage{longtable}
\usepackage{graphicx}
\usepackage{hyperref}
\usepackage{xcolor}
\usepackage{listings}
\usepackage{tcolorbox}
\usepackage{polya}

\hypersetup{
  colorlinks=true,
  linkcolor=polyablue,
  urlcolor=polyablue,
  citecolor=polyablue,
}

% ── Code listing style ────────────────────────────────────────────────
\lstdefinestyle{polya-python}{
  language=Python,
  basicstyle=\ttfamily\small,
  keywordstyle=\color{polyablue}\bfseries,
  commentstyle=\color{polyagray}\itshape,
  stringstyle=\color{polyagreen},
  numbers=left,
  numberstyle=\tiny\color{polyagray},
  numbersep=8pt,
  frame=single,
  framerule=0.4pt,
  rulecolor=\color{polyagray},
  breaklines=true,
  showstringspaces=false,
  tabsize=4,
}
\lstset{style=polya-python}
\title{Potluck Planning}
\author{uber-polya}
\date{2026-02-26}

\begin{document}

\maketitle
\tableofcontents
\newpage

% ── Phase A: Problem Formulation ─────────────────────────────────────
\section{Problem Formulation}

\begin{polyamodel}

\textbf{Problem Type}: Find \hfill
\textbf{Domain}: Integer Linear Programming

\textbf{Named Problem}: Assignment Problem

\end{polyamodel}

% ── Universe ──────────────────────────────────────────────────────────
\subsection{Universe}

\begin{itemize}
  \item $Assignment: \{Alice, Ben, Carmen, David, Elena, Frank, Grace, Hiro\}$
  \item $Category Counts: \{appetizer, salad, main, side, dessert\}$
\end{itemize}

% ── Variables ─────────────────────────────────────────────────────────
\subsection{Variables}

\begin{longtable}{lll}
\toprule
\textbf{Symbol} & \textbf{Meaning} & \textbf{Type / Range} \\
\midrule
\endhead
$x_{ij}$ & assignment decision & $\{0, 1\}$ \\
$n$ & number of binary variables = 40 & $\mathbb{{Z}}^+$ \\
\bottomrule
\end{longtable}

% ── Structure ─────────────────────────────────────────────────────────
\subsection{Structure}

Assign 8 guests to 5 dish categories (appetizer, salad, main, side, dessert) for a potluck dinner. Each guest rates their effort/skill for each category on a 1-5 scale (lower = easier for them). Every category must be covered by at least 1 guest and at most 2 guests. Minimize total effort across all assignments.

% ── Mapping ───────────────────────────────────────────────────────────
\subsection{Real-World Mapping}

\begin{longtable}{ll}
\toprule
\textbf{Real-World Concept} & \textbf{Mathematical Object} \\
\midrule
\endhead
assignment & $x: solution vector (assignment)$ \\
\bottomrule
\end{longtable}

% ── Constraints ───────────────────────────────────────────────────────
\subsection{Constraints}

\begin{enumerate}
  \item $\sum_j x_{ij} \le 2$
  \hfill \textit{(capacity limit)}
  \item $\sum_j x_{ij} \ge 1$
  \hfill \textit{(minimum requirement)}
  \item $x_{ij} \in \{0, 1\} \quad \forall\, i, j$
  \hfill \textit{(binary decision variables)}
\end{enumerate}

% ── Objective / Claim ─────────────────────────────────────────────────
\subsection{Claim}


% ── Approach ──────────────────────────────────────────────────────────
\subsection{Approach}

\begin{description}
  \item[Suggested approach:] ILP (PuLP/CBC, Branch \& Bound)
  \item[Complexity class:] 
  \item[Available tools:] ILP
\end{description}
% ── Phase B: Solution ────────────────────────────────────────────────
\section{Solution}

\begin{polyaresult}

\textbf{Answer}: 8-element assignment: Alice -> dessert, Ben -> salad, Carmen -> main, David -> side, Elena -> appetizer, Frank -> dessert, ... (2 more)


\textbf{Optimal}: Yes \hfill
\textbf{Feasible}: Yes
\textbf{Algorithm}: ILP (PuLP/CBC, Branch \& Bound) \quad $$

\textbf{Time}: 0.005929717910476029s

\textbf{Certificate}: Branch-and-bound solved to optimality.

\end{polyaresult}

% ── Solution Details ──────────────────────────────────────────────────
\subsection{Solution Details}

Assignment:
  Alice -> dessert
  Ben -> salad
  Carmen -> main
  David -> side
  Elena -> appetizer
  Frank -> dessert
  Grace -> salad
  Hiro -> appetizer

Method: Integer Linear Programming (ILP) via PuLP/CBC.

- 40 binary variables (8 guests x 5 categories)
- Constraints: one category per guest, min/max coverage per category
- Objective: minimize sum of effort scores across all guest-category assignments

% ── Verification ──────────────────────────────────────────────────────
\subsection{Verification}

\begin{longtable}{lcl}
\toprule
\textbf{Check} & \textbf{Status} & \textbf{Detail} \\
\midrule
\endhead
Feasibility & \passmark & All constraints satisfied \\
Optimality & \passmark & ILP (PuLP/CBC, Branch \& Bound) \\
\bottomrule
\end{longtable}

% ── Phase C: Interpretation ──────────────────────────────────────────
\section{Interpretation}

% ── The Question & The Answer ─────────────────────────────────────────
\subsection{The Question}
Assign 8 guests to 5 dish categories (appetizer, salad, main, side, dessert) for a potluck dinner.

\subsection{The Answer}

\begin{polyainsight}[title={Bottom Line}]
A optimal solution was found.
\end{polyainsight}

% ── What This Means ───────────────────────────────────────────────────
\subsection{What This Means}

- Optimal assignment of guests to dish categories minimizing total effort
- All 5 categories covered (some by 2 guests since 8 > 5)
- Independent constraint verification (all PASS)
- Comparison against theoretical bounds and random assignment

Integer Linear Programming (ILP) via PuLP/CBC.

- 40 binary variables (8 guests x 5 categories)
- Constraints: one category per guest, min/max coverage per category
- Objective: minimize sum of effort scores across all guest-category assignments

% ── Sensitivity Analysis ──────────────────────────────────────────────

% ── Figures ───────────────────────────────────────────────────────────

% ── Recommendations ───────────────────────────────────────────────────
\subsection{Recommendations}

\begin{enumerate}
  \item The solution is provably optimal; no further improvement is possible within this model.
\end{enumerate}

% ── Limitations ───────────────────────────────────────────────────────
\subsection{Limitations}

\begin{polyawarnbox}
\begin{itemize}
  \item Model assumes deterministic parameters; real-world variability not captured.
  \item Solution quality depends on the accuracy of input data.
\end{itemize}
\end{polyawarnbox}

% ── Appendix ─────────────────────────────────────────────────────────

\end{document}